\documentclass[11pt]{article}
\usepackage[margin=0.85in]{geometry}
\usepackage{amsmath,amssymb}
\usepackage[T1]{fontenc}

% EDIT YOUR ANSWERS HERE. Replace \text{TODO} with a math expression.
% Useful syntax: \land (and), \neg (not), \ge (>=), \cdot (multiply).
% Keep assignment substitutions unsimplified unless a ->> allows consequence.
\newcommand{\answerA}{(0 = (m - m) * (m * m)) \land (m * m = m * m) \land (m \geq 0)}
\newcommand{\answerB}{(y = (m - m) * (m * m)) \land (m * m = m * m) \land (m \geq 0)}
\newcommand{\answerC}{(y = (m - m) * x) \land (x = m * m) \land (m \geq 0)}
\newcommand{\answerD}{(y = (m - z) * x) \land (x = m * m) \land (z \geq 0)}
\newcommand{\answerE}{((y = (m - z) * x) \land (x = m * m) \land (z \geq 0)) \land (z > 0)}
\newcommand{\answerF}{((y + x) = (m - (z - 1)) * x) \land (x = m * m) \land ((z - 1) \geq 0)}
\newcommand{\answerG}{((y + x) = (m - z) * x) \land (x = m * m) \land (z \geq 0)}
\newcommand{\answerH}{(y = (m - z) * x) \land (x = m * m) \land (z \geq 0)}
\newcommand{\answerI}{(y = (m - z) * x) \land (x = m * m) \land (z \geq 0) \land \neg (z > 0)}

\newcommand{\assertion}[2]{%
  \noindent\makebox[1.8em][l]{\textsf{#1}}%
  $\{\!\{\;#2\;\}\!\}$\par\smallskip}
\newcommand{\code}[1]{\noindent\hspace*{1.8em}\texttt{#1}\par\smallskip}
\newcommand{\consequence}{\mathrel{\texttt{->>}}}

\begin{document}
\section*{2024 Midterm 1: Question 4 (20 points)}
Consider the following Dafny program that inefficiently calculates the cube
of a number:
\begin{verbatim}
method Cube(m:int) returns (y:int)
requires m > 0
ensures y == m*m*m
{
  y := 0;
  var x := m * m;
  var z := m;
  while (z > 0)
  {
    z := z - 1;
    y := y + x;
  }
}
\end{verbatim}
Fill in the annotations below to show that the Hoare triple given by the
outermost pre- and postconditions is valid. Be completely precise and
pedantic in applying Hoare rules: write assertions in exactly the form
given by the rules, rather than just logically equivalent ones. Working
backward from the end should require the rule of consequence only at the
places marked \texttt{->>}. Those implication steps may silently rely on
the usual rules of arithmetic.

\medskip
\noindent Edit the nine answer commands near the top of this file.
Labels A--I are worksheet labels, not part of the assertions.
The original decorated listing uses uppercase \texttt{Z} in its guard;
that spelling is preserved below. The method above uses lowercase \texttt{z}.
\medskip

\assertion{}{m > 0}
\noindent\hspace*{1.8em}$\consequence$\par\smallskip
\assertion{A}{\answerA}
\code{y := 0;}
\assertion{B}{\answerB}
\code{var x := m * m;}
\assertion{C}{\answerC}
\code{var z := m;}
\assertion{D}{\answerD}
\code{while (Z > 0) \{}
\assertion{E}{\answerE}
\noindent\hspace*{1.8em}$\consequence$\par\smallskip
\assertion{F}{\answerF}
\code{\quad z := z - 1;}
\assertion{G}{\answerG}
\code{\quad y := y + x;}
\assertion{H}{\answerH}
\code{\}}
\assertion{I}{\answerI}
\noindent\hspace*{1.8em}$\consequence$\par\smallskip
\assertion{}{y = m * m * m}

\end{document}
